\ifdefined\gkver\else\def\gkver{student}\fi
\documentclass[\gkver]{gaokaozhenti}
\begin{document}

\chapter{2010年四川理综}
%% paperType: 理综物理部分

\begin{enumerate}
\item
%% number: 14
%% paperName: 2010年四川理综第 14 题 6 分
%% typeId: 1
%% type: 单选题
%% chapter: 气体性质
%% point: 分子力
%% method: 无
%% score: 6
%% degree: 650
%% duplicateId: 0
%% body:
下列现象中不能说明分子间存在分子力的是\xzanswer{D}
\fourchoices[answer=D]
{两铅块能被压合在一起}
{钢绳不易被拉断}
{水不容易被压缩}
{空气容易被压缩}

%% answer: D
%% memo:
\memoanswer{
空气容易压缩是因为分子间距大，而水不容易压缩是因为分子间距小，轻微压缩都使分子力表现为斥力。ABC 说明存在分子力。
}

\item
%% number: 15
%% paperName: 2010年四川理综第 15 题 6 分
%% typeId: 1
%% type: 单选题
%% chapter: 光学
%% point: 光的电磁说
%% method: 无
%% score: 6
%% degree: 650
%% duplicateId: 0
%% body:
下列说法正确的是\xzanswer{D}
\fourchoices[answer=D]
{$\alpha$ 粒子大角度散射表明 $\alpha$ 粒子很难进入原子内部}
{氢原子跃迁发出的光从空气射入水时可能发生全反射}
{裂变反应有质量亏损，质量数不守恒}
{$\gamma$ 射线是一种波长很短的电磁波}

%% answer: D
%% memo:
\memoanswer{
$\alpha$ 粒子散射实验现象表明大多数 $\alpha$ 粒子不发生偏转，说明穿过了原子，少数 $\alpha$ 粒子发生偏转，说明无法穿过原子核，A 错误。光只有从光密介质进入光疏介质才可能发生全反射，空气和水比较水是光密介质，B 错误。裂变反应有质量亏损是由于核子的平均密度变化引起的，但核子的总数不变，即质量数守恒，C 错误。$\gamma$ 射线是频率很大、波长很短的电磁波，D 正确。
}

\item
%% number: 16
%% paperName: 2010年四川理综第 16 题 6 分
%% typeId: 1
%% type: 单选题
%% chapter: 机械振动
%% point: 振动图象
%% method: 无
%% score: 6
%% degree: 450
%% duplicateId: 0
%% body:
一列间谐横波沿直线由 A 向 B 传播，A、B 相距 $0.45\Um$，右图是 A 处质点的振动图像。当 A 处质点运动到波峰位置时，B 处质点刚好到达平衡位置且向 y 轴正方向运动，这列波的波速可能是\xzanswer{A}

\onepicture[w=6cm]{figs/16.png}

\fourchoices[answer=A]
{$4.5\Ums$}
{$3.0\Ums$}
{$1.5\Ums$}
{$0.7\Ums$}

%% answer: A
%% memo:
\memoanswer{
由振动图象可得振动周期 $T=0.4\Us$，波由 A 向 B 传播，A 处质点运动到波峰位置时，B 处质点刚好到达平衡位置且向 y 轴正方向运动，A、B 间的距离 $l=n\lambda+\frac{1}{4}\lambda$，则波长 $\lambda=\frac{4l}{4n+1}$，波速 $v=\frac{\lambda}{T}=\frac{4l}{T(4n+1)}=\frac{4\times0.45}{0.4\times(4n+1)}=\frac{9}{8n+2}$。当 $n=0$ 时，$v=4.5\Ums$；当 $n=1$ 时，$v=0.9\Ums$；当 $n=2$ 时，$v=0.5\Ums$ 等，A 正确。

\onepicture[w=3.2cm]{figs/16_an.png}
}

\item
%% number: 17
%% paperName: 2010年四川理综第 17 题 6 分
%% typeId: 1
%% type: 单选题
%% chapter: 曲线运动
%% point: 人造地球卫星
%% method: 无
%% score: 6
%% degree: 450
%% duplicateId: 0
%% body:
a 是地球赤道上一栋建筑，b 是在赤道平面内做匀速圆周运动、距地面 $9.6\times10^{6}\Um$ 的卫星，c 是地球同步卫星，某一时刻 b、c 刚好位于 a 的正上方（如图甲所示），经 $48\Uh$，a、b、c 的大致位置是图乙中的（取地球半径 $R=6.4\times10^{6}\Um$，地球表面重力加速度 $g=10\Umsq$，$\pi=\sqrt{10}$）\xzanswer{B}

\onepicture[num=甲, labelprefix={}, w=2.0cm]{figs/17a.png}

\fourchoices[answer=B, ispicture=true, h=2.2cm]
{figs/17b1.png}
{figs/17b2.png}
{figs/17b3.png}
{figs/17b4.png}

%% answer: B
%% memo:
\memoanswer{
b、c 都是地球卫星，共同遵循地球对它们的万有引力提供向心力，c 是地球同步卫星，c 在 a 的正上方，对 b 有 $G\frac{Mm}{(R+h)^{2}}=m\left(\frac{2\pi}{T}\right)^{2}(R+h)$，$G\frac{Mm}{R^{2}}=mg$，联立可得：$T=2\pi\sqrt{\frac{(R+h)^{3}}{gR^{2}}}=2\times10^{4}\Us$，经 $48\Uh$，b 转过的圈数 $n=\frac{48\times3600}{2\times10^{4}}=8.64$，选项 B 正确。
}

\item
%% number: 18
%% paperName: 2010年四川理综第 18 题 6 分
%% typeId: 1
%% type: 单选题
%% chapter: 光学
%% point: 光电效应
%% method: 无
%% score: 6
%% degree: 650
%% duplicateId: 0
%% body:
用波长为 $2.0\times10^{-7}\Um$ 的紫外线照射钨的表面，释放出来的光电子中最大的动能是 $4.7\times10^{-19}\UJ$。由此可知，钨的极限频率是（普朗克常量 $h=6.63\times10^{-34}\UJcdots$），光速 $c=3.0\times10^{8}\Ums$，结果取两位有效数字）\xzanswer{B}
\fourchoices[answer=B]
{$5.5\times10^{14}\UHz$}
{$7.9\times10^{14}\UHz$}
{$9.8\times10^{14}\UHz$}
{$1.2\times10^{15}\UHz$}

%% answer: B
%% memo:
\memoanswer{
根据光电效应方程 $E_{km}=h\nu-W$，在恰好发生光电效应时最大出动能为 0 有 $h\nu_{0}=W$，且 $c=\lambda\nu$，综合化简得 $\nu_{0}=\frac{h\frac{c}{\lambda}-E_{km}}{h}=\frac{c}{\lambda}-\frac{E_{km}}{h}=\frac{3\times10^{8}}{2\times10^{-7}}-\frac{4.7\times10^{-19}}{6.63\times10^{-34}}=7.9\times10^{14}\UHz$，B 正确。
}

\item
%% number: 19
%% paperName: 2010年四川理综第 19 题 6 分
%% typeId: 1
%% type: 单选题
%% chapter: 电磁感应
%% point: 电磁感应中图象
%% method: 无
%% score: 6
%% degree: 450
%% duplicateId: 0
%% body:
图甲所示电路中，$A_{1}$、$A_{2}$、$A_{3}$ 为相同的电流表，$C$ 为电容器，电阻 $R_{1}$、$R_{2}$、$R_{3}$ 的阻值相同，线圈 $L$ 的电阻不计。在某段时间内理想变压器原线圈内磁场的变化如图乙所示，则在 $t_{2}\sim t_{3}$ 时间内\xzanswer{C}

\twopicture[numA=甲, numB=乙, labelprefix={}, wA=5.6cm, wB=3.8cm]
{figs/19a.png}
{figs/19b.png}

\fourchoices[answer=C]
{电流表 $A_{1}$ 的示数比 $A_{2}$ 的小}
{电流表 $A_{2}$ 的示数比 $A_{1}$ 的小}
{电流表 $A_{1}$ 和 $A_{2}$ 的示数相同}
{电流表的示数都不为零}

%% answer: C
%% memo:
\memoanswer{
由 $B-t$ 图象知在 $t_{1}\sim t_{2}$ 时间内，原线圈中的磁场先负向减小后正向增大，则副线圈中的磁通量是均匀变化的，根据法拉第电磁感应定律在副线圈中产生的感应电流大小不变，再根据楞次定律可判断负向减小时正向增大时的感应电流的方向相同，则在 $t_{1}\sim t_{2}$ 时间内副线圈中感应电流为恒定电流，所以 $A_{1}$、$A_{2}$ 的示数相同，$A_{3}$ 的示数为 0，C 正确。
}

\item
%% number: 20
%% paperName: 2010年四川理综第 20 题 6 分
%% typeId: 2
%% type: 多选题
%% chapter: 电磁感应
%% point: 双杆问题
%% method: 无
%% score: 6
%% degree: 450
%% duplicateId: 0
%% body:
如图所示，电阻不计的平行金属导轨固定在一绝缘斜面上，两相同的金属导体棒 a、b 垂直于导轨静止放置，且与导轨接触良好，匀强磁场垂直穿过导轨平面。现用一平行于导轨的恒力 $F$ 作用在 a 的中点，使其向上运动。若 b 始终保持静止，则它所受摩擦力可能\xzanswer{AB}

\onepicture[align=right, w=4cm]{figs/20.png}

\fourchoices[answer=AB]
{变为 0}
{先减小后不变}
{等于 $F$}
{先增大再减小}

%% answer: AB
%% memo:
\memoanswer{
对 a 棒受到合力为 $F_{\text{合}}=F-f-mg\sin\theta-Bl^{2}v^{2}/R$ 说明 a 做加速度减小的加速运动，当加速度为 0 后匀速运动，所以 a 受安培力先增大后不变。

如果 $F=F_{f}+2mg\sin\theta$，则安培力为 $mg\sin\theta$，则 b 受的摩擦力最后为 0；

$F<F_{f}+2mg\sin\theta$，则安培力小于 $mg\sin\theta$，则 b 受的摩擦力一直减小，最后不变，B 正确；

如果 $F_{f}+2mg\sin\theta<F<F_{f}+3mg\sin\theta$，则安培力大于 $mg\sin\theta$ 或小于 $2mg\sin\theta$，则 b 受的摩擦力先减小后增大，最后不变。可以看出 b 受的摩擦力先变化后不变，C、D 错误。
}

\item
%% number: 21
%% paperName: 2010年四川理综第 21 题 6 分
%% typeId: 2
%% type: 多选题
%% chapter: 电场
%% point: 带电粒子的直线运动
%% method: 无
%% score: 6
%% degree: 450
%% duplicateId: 0
%% body:
如图所示，圆弧虚线表示正点电荷电场的等势面，相邻两等势面间的电势差相等。光滑绝缘直杆沿电场方向水平放置并固定不动，杆上套有一带正电的小滑块（可视为质点），滑块通过绝缘轻弹簧与固定点 O 相连，并以某一初速度从 M 点运动到 N 点，$OM<ON$。若滑块在 M、N 时弹簧的弹力大小相等，弹簧始终在弹性限度内，则\xzanswer{AC}

\onepicture[align=right, w=4.5cm]{figs/21.png}

\fourchoices[answer=AC]
{滑块从 M 到 N 的过程中，速度可能一直增大}
{滑块从位置 1 到 2 的过程中，电场力做的功比从位置 3 到 4 的小}
{在 M、N 之间的范围内，可能存在滑块速度相同的两个位置}
{在 M、N 之间可能存在只由电场力确定滑块加速度大小的三个位置}

%% answer: AC
%% memo:
\memoanswer{
在 N 点如果电场力不小于弹簧弹力的分力，则滑块一直加速，A 正确。在 N 点如果电场力小于弹簧弹力的分力，则滑块先加速后减速，就可能有两个位置的速度相同，C 正确。1、2 与 3、4 间的电势差相等，电场力做功相等，B 错误。由于 M 点和 N 点弹簧的长度不同但弹力相等，说明 N 点时弹簧是压缩的，在弹簧与水平杆垂直和弹簧恢复原长的两个位置滑块的加速度只由电场力决定，D 错误。
}

\item
%% number: 22
%% paperName: 2010年四川理综第 22 题 13 分
%% typeId: 4
%% type: 实验题
%% chapter: 电路
%% point: 电路实验
%% method: 无
%% score: 13
%% degree: 450
%% duplicateId: 0
%% body:
（1）用多用电表探测图甲所示黑箱发现：用直流电压挡测量，E、G 两点间和 F、G 两点间均有电压，E、F 两点间无电压；用欧姆挡测量，黑表笔（与电表内部电源的正极相连）接 E 点，红表笔（与电表内部电源的负极相连）接 F 点，阻值很小，但反接阻值很大。那么，该黑箱内元件的接法可能是图乙中\tkanswer[1.2cm]{B}。

\onepicture[num=甲, labelprefix={}, w=2.3cm]{figs/22a.png}

\fourchoices[answer=B, ispicture=true, h=1.8cm]
{figs/22b1.png}
{figs/22b2.png}
{figs/22b3.png}
{figs/22b4.png}

（2）在物理兴趣小组活动中，一同学利用下列器材设计并完成了“探究导体阻值与长度的关系”的实验。

电压表 $V_{1}$\quad 量程 $3\UV$\quad 内阻约为 $900\UO$

电压表 $V_{2}$\quad 量程 $10\UV$\quad 内阻约为 $3\UkO$

电压表 A\quad 量程 $60\UmA$\quad 内阻约为 $5\UO$

电源 $E_{1}$\quad 电动势 $1.5\UV$\quad 内阻约为 $0.2\UO$

电源 $E_{2}$\quad 电动势 $4.5\UV$\quad 内阻约为 $0.4\UO$

滑动变阻器（最大阻值为 $10\UO$）、粗细均匀的同种电阻丝、开关、导线和刻度尺。其主要实验步骤如下：

\begin{steps}
	\item
	选取图中器材，按示意图连接电路；
	\item
	用伏安法测定电阻丝的阻值 $R$；
	\item
	用刻度尺测出电阻丝的长度 $L$；
	\item
	依次减小电阻丝的长度，保持电路其他部分不变，重复步骤②、③；
	\item
	处理数据，根据下列测量结果，找出电阻丝阻值与长度的关系。
\end{steps}

\begin{table}[htbp]
\centering
\begin{tblr}{|c|c|c|c|c|c|}
\hline
$L$（$\Um$） & 0.9956 & 0.8049 & 0.5981 & 0.4021 & 0.1958 \\
\hline
$R$（$\UO$） & 104.8 & 85.3 & 65.2 & 46.6 & 27.1 \\
\hline
\end{tblr}
\end{table}

为使实验尽可能准确，请你对上述步骤中画线处加以改进。

（Ⅰ）\tkanswer[3.5cm]{②电压改选 $E_{2}$}

（Ⅱ）\tkanswer[3.5cm]{判断电流表的内外接法，作出相应调整}

\onepicture[align=right, w=3.2cm]{figs/22c.png}

%% answer: 
%% memo:
\memoanswer{
①红表笔接 F，电阻很小，此时二极管导通，电源电流从黑表笔流出通过二极管从红表笔流进，电流方向 E 到 F，只有 B 正确。

②（Ⅰ）根据电动势的大小和电压表的量程不能选择 $10\UV$，电路图选择 $V_{1}$ 当然正确。为减小测量误差要求电压表的指针偏转过半，故 $1.5\UV$ 的电动势太小，电源 $E_{1}$ 应该选为 $E_{2}$。

（Ⅱ）由于 $\sqrt{R_{V1}R_{A}}=\sqrt{900\times5}$ 约 60 多欧姆，第一组测量值为 $104.8\UO$ 为大电阻用内接伏安法，第五组测量值为 $27.1\UO$，为小电阻用外接伏安法，这说明外接法先用试触法判断调整后再测量。
}

\item
%% number: 22
%% paperName: 2010年四川理综第 22 题 4 分
%% typeId: 4
%% type: 实验题
%% chapter: 直线运动
%% point: 打点计时器公式
%% method: 无
%% score: 4
%% degree: 650
%% duplicateId: 0
%% body:
有 4 条用打点计时器（所用交流电频率为 $50\UHz$）打出的纸带 A、B、C、D，其中一条是做“验证机械能守恒定律”实验时打出的。为找出该纸带，某同学在每条纸带上取了点迹清晰的、连续的 4 个点，用刻度尺测出相邻两个点间距离依次为 $S_{1}$、$S_{2}$、$S_{3}$。请你根据下列 $S_{1}$、$S_{2}$、$S_{3}$ 的测量结果确定该纸带为\tkanswer[1.5cm]{C}。（已知当地的重力加速度为 $9.791\Umsq$）

（A）$61.0\Umm$、$65.8\Umm$、$70.7\Umm$\quad（B）$41.2\Umm$、$45.1\Umm$、$53.0\Umm$\quad（C）$49.36\Umm$、$53.5\Umm$、$57.3\Umm$\quad（D）$60.5\Umm$、$61.0\Umm$、$60.6\Umm$

%% answer: 
%% memo:
\memoanswer{
验证机械能守恒定律采用重锤的自由落体实现，所以相邻的 $0.02\Us$ 时间内的位移增大量为 $\Delta S=gT^{2}=9.791\times0.02^{2}\Um=3.9\Umm$，答案 C。
}

\item
%% number: 23
%% paperName: 2010年四川理综第 23 题 16 分
%% typeId: 6
%% type: 计算题
%% chapter: 运动和力
%% point: 牛顿定律的应用
%% method: 无
%% score: 16
%% degree: 650
%% duplicateId: 0
%% body:
质量为 $M$ 的拖拉机拉着耙来耙地，由静止开始做匀加速直线运动，在时间 $t$ 内前进的距离为 $s$。耙地时，拖拉机受到的牵引力恒为 $F$，受到地面的阻力为自重的 $k$ 倍，耙所受阻力恒定，连接杆质量不计且与水平面的夹角 $\theta$ 保持不变。求：
\begin{enumerate}
	\item
	拖拉机的加速度大小。
	\item
	拖拉机对连接杆的拉力大小。
	\item
	时间 $t$ 内拖拉机对耙做的功。
\end{enumerate}
\onepicture[align=right, w=5.5cm]{figs/23.png}

%% answer: 
\jdanswer{
\begin{enumerate}
	\item
	$a=\frac{2s}{t^{2}}$
	\item
	$T=\frac{1}{\cos\theta}\left[F-M\left(kg+\frac{2s}{t^{2}}\right)\right]$
	\item
	$W=\left[F-M\left(kg+\frac{2s}{t^{2}}\right)\right]s$
\end{enumerate}
}
%% memo:
\memoanswer{
（1）拖拉机在时间 $t$ 内匀加速前进 $s$，根据位移公式

$s=\frac{1}{2}at^{2}$ ①

变形得 $a=\frac{2s}{t^{2}}$ ②

（2）对拖拉机受到牵引力、支持力、重力、地面阻力和连杆拉力 $T$，根据牛顿第二定律

$Ma=F-kMg-T\cos\theta$ ③

联立②③变形得 $T=\frac{1}{\cos\theta}\left[F-M\left(kg+\frac{2s}{t^{2}}\right)\right]$ ④

根据牛顿第三定律连杆对耙的反作用力为

$T'=T=\frac{1}{\cos\theta}\left[F-M\left(kg+\frac{2s}{t^{2}}\right)\right]$ ⑤

拖拉机对耙做的功：$W=T's\cos\theta$ ⑥

联立④⑤解得 $W=\left[F-M\left(kg+\frac{2s}{t^{2}}\right)\right]s$ ⑦
}

\item
%% number: 24
%% paperName: 2010年四川理综第 24 题 19 分
%% typeId: 6
%% type: 计算题
%% chapter: 磁场
%% point: 带电粒子在复合场中的运动
%% method: 无
%% score: 19
%% degree: 450
%% duplicateId: 0
%% body:
如图所示，电源电动势 $E_{0}=15\UV$，内阻 $r_{0}=1\UO$，电阻 $R_{1}=30\UO$，$R_{2}=60\UO$。间距 $d=0.2\Um$ 的两平行金属板水平放置，板间分布有垂直于纸面向里、磁感应强度 $B=1\UT$ 的匀强磁场。闭合开关 S，板间电场视为匀强电场，将一带正电的小球以初速度 $v=0.1\Ums$ 沿两板间中线水平射入板间。设滑动变阻器接入电路的阻值为 $R_{x}$，忽略空气对小球的作用，取 $g=10\Umsq$。
\begin{enumerate}
	\item
	当 $R_{x}=29\UO$ 时，电阻 $R_{2}$ 消耗的电功率是多大？
	\item
	若小球进入板间做匀速圆周运动并与板相碰，碰时速度与初速度的夹角为 $60^{\circ}$，则 $R_{x}$ 是多少？
\end{enumerate}
\onepicture[align=right, w=5cm]{figs/24.png}

%% answer: 
\jdanswer{
\begin{enumerate}
	\item
	$0.6\UW$
	\item
	$54\UO$
\end{enumerate}
}
%% memo:
\memoanswer{
（1）闭合电路的外电阻为

$R=R_{x}+\frac{R_{1}R_{2}}{R_{1}+R_{2}}=29+\frac{30\times60}{30+60}=49\UO$ ①

根据闭合电路的欧姆定律

$I=\frac{E}{R+r}=\frac{15}{49+1}\UA=0.3\UA$ ②

$R_{2}$ 两端的电压为

$U_{2}=E-I(R_{x}+r)=15-0.3\times30=6\UV$ ③

$R_{2}$ 消耗的功率为

$P_{2}=\frac{U_{2}^{2}}{R_{2}}=\frac{6^{2}}{60}\UW=0.6\UW$ ④

（2）小球进入电磁场做匀速圆周运动，说明重力和电场力等大反向，洛伦兹力提供向心力，根据牛顿第二定律

$Bqv=m\frac{v^{2}}{R}$ ⑤

$\frac{U_{2}}{d}q=mg$ ⑥

联立⑤⑥化简得

$U_{2}=\frac{BRdg}{v}$ ⑦

小球做匀速圆周运动的初末速的夹角等于圆心角为 $60^{\circ}$，根据几何关系得

$R=d$ ⑧

联立⑦⑧带入数据

$U_{2}=\frac{Bd^{2}g}{v}=\frac{1\times0.04\times10}{0.1}\UV=4\UV$

干路电流为 $I=\frac{U_{2}}{R_{12}}=\frac{4}{20}\UA=0.2\UA$ ⑨

$R_{x}=\frac{E-U_{2}}{I}-r=\frac{15-4}{0.2}-1=54\UO$ ⑩
}

\item
%% number: 25
%% paperName: 2010年四川理综第 25 题 20 分
%% typeId: 6
%% type: 计算题
%% chapter: 电场
%% point: 带电粒子的偏转
%% method: 无
%% score: 20
%% degree: 250
%% duplicateId: 0
%% body:
如图所示，空间有场强 $E=0.5\UNC$ 的竖直向下的匀强电场，长 $l=0.3\sqrt{3}\Um$ 的不可伸长的轻绳一端固定于 O 点，另一端系一质量 $m=0.01\Ukg$ 的不带电小球 A，拉起小球至绳水平后，无初速释放。另一电荷量 $q=+0.1\UC$、质量与 A 相同的小球 P，以速度 $v_{0}=3\sqrt{3}\Ums$ 水平抛出，经时间 $t=0.2\Us$ 与小球 C 在 D 点下方一足够大的平板相遇。不计空气阻力，小球均可视为质点，取 $g=10\Umsq$。
\begin{enumerate}
	\item
	求碰撞前瞬间小球 P 的速度。
	\item
	若小球 C 经过路 $s=0.09\Um$ 到达平板，此时速度恰好为 0，求所加的恒力。
	\item
	若施加恒力后，保持平板垂直于纸面且与水平面的夹角不变，在 D 点下方面任意改变平板位置，小球 C 均能与平板正碰，求出所有满足条件的恒力。
\end{enumerate}
\onepicture[align=right, w=5.5cm]{figs/25.png}

%% answer: 
\jdanswer{
\begin{enumerate}
	\item
	$v_{1}=6\Ums$，$\beta=30^{\circ}$
	\item
	$\alpha=30^{\circ}$，$F=\frac{\sqrt{3}}{4}\UN$
	\item
	$F=\frac{\sqrt{3}}{8\sin(30^{\circ}+\alpha)}\UN$
\end{enumerate}
}
%% memo:
\memoanswer{
（1）设 P 的加速度为 $a_{0}$、到 D 点的竖直速度为 $v_{y}$，合速度大小为 $v_{1}$，与水平方向的夹角为 $\beta$，有

$mg+qE=ma_{0}$

$v_{y}=a_{0}t$

$v_{1}^{2}=v_{0}^{2}+v_{y}^{2}$

$\tan\beta=\frac{v_{y}}{v_{0}}$

联立上述方程，代入数据解得：$v_{1}=6\Ums$，$\beta=30^{\circ}$。

\onepicture[w=4.5cm]{figs/25_an1.png}

（2）设 A 碰前速度为 $v_{2}$，此时轻绳与竖直线的夹角为 $\beta$，由动能定理得

$mgl\cos\beta=\frac{1}{2}mv_{2}^{2}$

设 A、P 碰撞后小球 C 的速度为 $v$，由动量守恒定律得

$mv_{1}-mv_{2}=2mv$

小球 C 到达平板时速度为 0，应做匀减速直线运动，设加速度的大小为 $a$ 有

$v^{2}=2as$

设恒力大小为 $F$，与竖直方向的夹角为 $\alpha$，如图，根据牛顿第二定律，得

$F\cos(90^{\circ}-\alpha-\beta)-2mg\sin\beta-qE\sin\beta=2ma$

$F\cos(90^{\circ}-\alpha-\beta)-2mg\cos\beta-qE\cos\beta=0$

\onepicture[w=4cm]{figs/25_an2.png}

代入相关数据解得：$F=\frac{\sqrt{3}}{4}\UN$，$\alpha=30^{\circ}$。

（3）由于平板可距 D 点无限远，小球 C 必做匀速或匀加速直线运动，恒力 $F_{1}$ 的方向可从竖直方向顺时针转向无限接近速度方向，设恒力与速度方向夹角为 $\theta$，有

$0\leq\theta<(90^{\circ}+30^{\circ})=120^{\circ}$

在垂直于速度方向上，有

$F_{1}\cos(\beta-\theta)=(2mg+qE)\cos\beta$

则 $F_{1}$ 的大小满足条件为

$F_{1}=\frac{\sqrt{3}}{8\cos(30^{\circ}-\theta)}\UN$（式中 $0\leq\theta<120^{\circ}$）

\onepicture[w=6cm]{figs/25_an3.png}
}

\end{enumerate}

\end{document}
